\section{Eleven unrestricted integers: selective coprime fusion}
\label{sec:integer-fusion}

The principal new equation has \textbf{eleven unrestricted integer unknowns,
exact degree $688$, and $126$ live binary arithmetic operations}.
Its direct sibling has degree $728$ and $122$ operations.
The preceding eleven-integer endpoint had degree $2240$ and $206$ operations.
A more aggressive fusion gives a lower-cost alternative of degree $1632$
and $114$ operations. These are comparisons of complete displayed
circuits, not global optimality claims. The unbounded native/refutation
calibration remains inherited; no new formal kernel certificate is claimed.

\subsection{The complete polynomial}
The independent coordinates are
\[
 (j,C,z,h,r,v,w,T,c,n,s)\in\mathbb Z^{11}.
\]
Use the six fixed integers $\beta,K_0,M_C,M_F,A_c,B_c$ specified by the
retained source-34 coefficient recipe, with $M_F$ already shifted.
Here $A_c=2^{1113}$ and $B_c=2^{666104}$; neither is a variable.

All subsequent symbols are polynomial abbreviations:
\begin{align*}
 J&=6j+1,&q&=\beta J+1,& W&=12w+4,\\
 X&=qW,&Y&=(2v+1)q^3,& E&=XY,\\
 a&=E+Y,& H&=4a+3,&\Delta&=a^2+H,\\
 F&=(K_0+W)C-z(q-1),\\
 R&=(q(q-F)-C)(q^2-1)+(M_C+qM_F)J,\\
 k&=R+1+hE,&\eta&=c-Yk,\\
 D_{\rm err}&=(A_cF-q)^2+(B_cC-q)^2,\\
 B&=k^2(q^2-D_{\rm err})-4q^2\eta^2,\\
 L&=XY^2k,&D_t&=1+L(L+k),\\
 D&=3B-4,&D_g&=1+4DD_t s^2,&P&=D_tD_g,\\
 I&=2src^2,&D_f&=1+\Delta I^2,&Q&=\Delta^2I^2,\\
 \delta&=Q-1,&V&=D_f(c-R+4IT)-c,\\
 C_y&=QV^2-1,&D_a&=\delta C_y,\\
 K&=X(3c-2X)+2,&S&=IH^5,&z_c&=HT+K,\\
 A_*&=(z_c^2+3H^2)(z_c^3-19H^3),\\
 Z&=S\bigl(n+(2Q-1)V\bigr)-A_*,\\
 u&=S^2D_f,&b&=4S^2D_a,&e&=16S^2P.
\end{align*}
The promised single polynomial equation is
\begin{equation}\label{eq:integer-fusion}
 \boxed{\ \bigl[(Z^2-u-b-e)^2-4(ub+ue+be)\bigr]^2
                 -64ubeZ^2=0.\ }
\end{equation}
There are no radicals, divisions, variable exponents or implicit witnesses
in this display. Removing $(2Q-1)V$ from $Z$ gives the direct $728/122$
variant by an integral translation of $n$.

\subsection{Why the fusion is valid}
The polynomial in~\eqref{eq:integer-fusion} is the product of the eight
independent conjugates of $Z-\sqrt u-\sqrt b-\sqrt e$.
The former five-root construction is reduced in two different ways:
the main root is recovered arithmetically rather than included in the norm,
and two other square conditions are fused using
\[
 \gcd(D_t,D_g)=\gcd(D_t,1+4DD_t s^2)=1.
\]
Thus $D_tD_g$ is a positive square exactly when both coprime positive
factors are squares. This is not an assumption that arbitrary products of
squares can be reversed.

For every integer assignment $H\ne0$, $\Delta>1$, $D_f,D_t\ge1$.
If $I=0$, the equation specializes to $A_*^8\ne0$, because
\[
 A_*=H^5\pi(T+K/H),\qquad \pi(x)=(x^2+3)(x^3-19)
\]
has no rational zero. Hence $I,s,r,c\ne0$ at any zero.
Then $|I|\ge c^2$, $D_f>|c|+1$, $|V|\ge1$, and $D_a>0$.
If $D_g<0$, its fused radical would be the only imaginary term in a
vanishing factor, an impossibility. Therefore $D_g>0$, and $B\ge2$.
The centered budget supplies the common strict signs of $q,F,C,q-F-C$
and $|c/k-Y|<1/2$.

All three radicands $D_f,D_a,P$ are odd positive integers. In a rational
signed sum with weights $1,2,4$, each irrational squareclass has a unique
least $2$-adic weight and cannot cancel. Thus all three roots are integers.
Coprimality recovers separate square roots of $D_t,D_g$.
The auxiliary quotient $y=\sqrt{D_a}/\delta$ is integral by the retained
adjacent-Pell descent. The factor identity now implies $S\mid A_*$.
Monicity and prime-valuation comparison give $H\mid K$. Since
\[
 K=HXc-2\{X(X+2ac)-1\},\qquad H\text{ odd},
\]
the signed quadratic main projector is recovered as well.

The missing main Pell norm must not be assumed in this argument.
From the strong and auxiliary norms, absolute-value Pell folding gives
$|c|=\psi_{|a+2|}(p)$ for some $p\ge1$ before projection.
The retained centered signed-rank proof then yields $p=R=2t+1$, $t\ge24$,
and the requisite ratio bound. The signed whole-product estimate yields
$\operatorname{sign}(c)X=2^R$; the original centered half-binomial argument
restores positive orientation and the native interval. Fresh private
witnesses then restore the exact old signed kernel and its finite
syntactic refutation endpoint. The complete proof records each dependency
and its order in \texttt{integer/PROOF.md}.

Conversely, a finite refutation supplies the same canonical outer tuple,
with $B>64$ and $D_t=\tau^2>0$. Since $(3B-4)D_t\equiv2\pmod3$,
ordinary Pell existence and a squared unit supply $s,g>0$ with
$D_g=g^2$. The retained full-residue interpolation theorem then chooses
one fixed strong pair $I,f$ before any residue for $T$ is requested.
The intersective polynomial $\pi$ has a root modulo every integer;
choose $T$ so that $I\mid\pi(T+K/H)$, and finally set
\[
 n=A_*/S+f+2\delta y+4\tau g-(2Q-1)V.
\]
This makes a conjugate factor zero. In particular, the strong witness is
not varied after selecting the intersective residue.

\subsection{Exact degree and independently checked arithmetic}
The relevant degrees are
\[
 \deg(S,D_f,Q,V,D_a,D_t,D_g,P)=(34,20,32,25,114,34,60,94).
\]
Thus $\deg(u,b,e)=(88,182,162)$. Put $C_*=S(2Q-1)V$ and
$s_*=Sn-A_*$. The exact identity
\[
 b-C_*^2=S^2(-V^2-4Q+4)
\]
gives $\deg(Z^2-b)\le126$. Rewriting the norm with
$\omega=Z^2-b$, its inner bracket has unique leading term $-4be$
of degree $344$; the remaining outer correction has degree $614$.
Hence the output has degree $688$ and leading form $16b_0^2e_0^2$.
For the authentic fixed tuple its coefficient
\[
 [j^{376}w^{84}v^{96}h^{12}s^{32}r^{28}c^{56}T^4]E_8
       =2^{780}3^{462}\beta^{376}>0
\]
proves that this bound is attained. No generic-coefficient assumption or
random sampling is used.

The portable checker authenticates all six numeral payloads, every live
arithmetic register, all $35$ coefficients of the free three-root norm,
all degree-leading inductions and the centered cancellation cut.
Both ordinary and optimized Python runs pass. Counting every live binary
addition, subtraction and multiplication, including squares and scalar
products, gives exactly $126$ operations; leaves and reuse are free.

\subsection{A lower-operation alternative}
Replace $I$ by
\[
 I=2src^2D_tD_gD_m(a+2),\qquad D_m=1+\Delta c^2,
\]
recompute the dependent definitions, and set $Z=Sn-A_*$. The equation
\[
 \boxed{\ \delta Z^2-S^2D_fD_tD_gC_y=0\ }
\]
has degree $1632$ and $114$ operations, still in eleven integers.
It is a polynomial despite its derivation from a rational auxiliary root.
Now $D_f,D_t,D_g,D_a$ are pairwise coprime: the new factors in $I$
give the first four coprimalities, while
\[
 D_a\equiv(a+2)^2D_m\pmod {D_f}
\]
gives the last one. Squaring $\delta Z/S$ recovers the product of these
four positive coprime integers; each is a square. The same adjacent-Pell
descent and scalar projection finish soundness. All additional factors
in $I$ are fixed before applying the same full-residue interpolation
lemma, so completeness introduces no stronger unproved interpolation
assumption. The detailed proof and complete circuit accompany the package.
